\subsubsection{Cache}

% ── DTO IN ───────────────────────────────────────────────────────────────────

% infrastructure\http\dtos\in\api-common.dto.ts
\addAtr{+}{collection: CollectionMetaDto}{Metadati di una collezione paginata ricevuta dall'API}
\makeCl{CollectionMetaDto}{Wrapper per i metadati di paginazione restituiti dall'API esterna}

% infrastructure\http\dtos\in\api-datapoint.dto.ts
\addAtr{+}{attributes: DatapointAttributesDto}{Attributi del datapoint restituiti dall'API}
\addAtr{+}{meta: DatapointMetaDto}{Metadati specifici del datapoint restituiti dall'API}
\addAtr{+}{meta: ApiMetaDto}{Metadati generici della risposta API}
\addAtr{+}{data: ApiDatapointDto[]}{Lista di datapoint grezzi restituiti dall'API}
\makeCl{DatapointAttributesDto}{DTO che mappa gli attributi di un datapoint ricevuto dall'API esterna}

% infrastructure\http\dtos\in\api-device.dto.ts
\addAtr{+}{attributes: DeviceAttributesDto}{Attributi del dispositivo restituiti dall'API}
\addAtr{+}{meta: DeviceMetaDto}{Metadati specifici del dispositivo restituiti dall'API}
\addAtr{+}{meta: ApiMetaDto}{Metadati generici della risposta API}
\addAtr{+}{data: ApiDeviceDto[]}{Lista di dispositivi grezzi restituiti dall'API}
\makeCl{DeviceAttributesDto}{DTO che mappa gli attributi di un dispositivo ricevuto dall'API esterna}

% infrastructure\http\dtos\in\api-plant-attributes.dto.ts
\makeCl{ApiPlantAttributesDto}{DTO che mappa gli attributi principali di un impianto ricevuto dall'API esterna}

% infrastructure\http\dtos\in\api-plant.dto.ts
\addAtr{+}{attributes: ApiPlantAttributesDto}{Attributi dell'impianto restituiti dall'API}
\addAtr{+}{meta: ApiPlantMetaDto}{Metadati specifici dell'impianto restituiti dall'API}
\addAtr{+}{meta: ApiMetaDto}{Metadati generici della risposta API}
\addAtr{+}{data: ApiRoomDto[]}{Lista di stanze grezze associate all'impianto}
\makeCl{ApiPlantMetaDto}{DTO che mappa i metadati di un impianto ricevuto dall'API esterna}

% infrastructure\http\dtos\in\plant-seek.dto.ts
\makeItf{PlantSeekResponseDto}{DTO di risposta per la ricerca di un impianto tramite ID}

% infrastructure\http\dtos\in\subNotification.dto.ts
\makeItf{NotificationLinkDto}{DTO che rappresenta un link incluso nel payload di una notifica esterna}

% ── MODEL ────────────────────────────────────────────────────────────────────

% [Plant e Room sono modelli di dominio impliciti, usati come tipi nelle porte e negli adapter]

% ── CONTROLLER ───────────────────────────────────────────────────────────────

% adapters\in\http\http-cache.controller.ts
\addAtr{-}{updateCacheUseCase: UpdateCacheUseCase}{Caso d'uso per aggiornare la cache di impianti specifici}
\addAtr{-}{webhookQueue: Promise<void>}{Coda sequenziale per serializzare i webhook in ingresso}
\addMet{+}{updateCache(payload: subNotificationPayloadDto)}{Riceve il webhook e avvia l'aggiornamento della cache per gli impianti indicati}
\makeCl{HttpCacheController}{Controller HTTP che gestisce i webhook in ingresso e delega l'aggiornamento della cache}

% adapters\in\event\event-cache.controller.ts
\addAtr{-}{cache: UpdateCacheAllPlantsUseCase}{Caso d'uso per aggiornare la cache di tutti gli impianti}
\makeCl{EventCacheController}{Controller eventi che ascolta gli eventi interni e avvia l'aggiornamento della cache per tutti gli impianti}

% ── CMD ──────────────────────────────────────────────────────────────────────

% application\commands\fetch-new-cache.command.ts
\makeItf{FetchNewCacheCmd}{Comando che trasporta token e plantId per recuperare la struttura aggiornata di un impianto}

% application\commands\get-valid-cache.command.ts
\makeItf{GetValidCacheCmd}{Comando che trasporta gli ID degli impianti per cui aggiornare la cache}

% ── USE CASE ─────────────────────────────────────────────────────────────────

% application\ports\in\update-cache.usecase.ts
\addMet{+}{updateCache(cmd: GetValidCacheCmd) : Promise$<$boolean$>$}{Aggiorna la cache per gli impianti specificati nel comando}
\makeItf{UpdateCacheUseCase}{Porta in ingresso per aggiornare la cache di uno o più impianti specifici}

% application\ports\in\update-cache-all-plants.usecase.ts
\addMet{+}{updateAllCache() : Promise$<$boolean$>$}{Avvia l'aggiornamento della cache per tutti gli impianti disponibili}
\makeItf{UpdateCacheAllPlantsUseCase}{Porta in ingresso per aggiornare la cache di tutti gli impianti}

% ── SERVICE ──────────────────────────────────────────────────────────────────

% application\services\sync-cache.service.ts
\addAtr{-}{fetchCachePort: FetchNewCachePort}{Porta per recuperare la struttura aggiornata di un impianto}
\addAtr{-}{writeStructurePort: WriteCachePort}{Porta per persistere la struttura di un impianto in cache}
\addAtr{-}{getAllPlantIdsPort: GetAllPlantIdsPort}{Porta per ottenere la lista di tutti gli ID impianto}
\addAtr{-}{emitter: EventEmitter2}{Emettitore di eventi interni al termine dell'aggiornamento}
\addMet{+}{updateCache(cmd: GetValidCacheCmd) : Promise$<$boolean$>$}{Recupera e persiste la struttura degli impianti indicati nel comando}
\addMet{+}{updateAllCache() : Promise$<$boolean$>$}{Recupera tutti gli ID impianto e aggiorna la cache per ciascuno}
\makeCl{SyncCacheService}{Servizio applicativo che orchestra il recupero e la persistenza della struttura degli impianti in cache}

% ── PORT ─────────────────────────────────────────────────────────────────────

% application\ports\out\fetch-new-cache.port.ts
\addMet{+}{fetch(cmd: FetchNewCacheCmd) : Promise$<$Plant$>$}{Recupera la struttura aggiornata di un impianto dato il comando}
\makeItf{FetchNewCachePort}{Porta di uscita per recuperare la struttura aggiornata di un impianto dall'esterno}

% application\ports\out\get-all-plantids.port.ts
\addMet{+}{getAllPlantIds() : Promise$<$string[]$>$}{Restituisce la lista di tutti gli ID impianto disponibili}
\makeItf{GetAllPlantIdsPort}{Porta di uscita per ottenere la lista completa degli ID impianto}

% application\ports\out\write-cache.port.ts
\addMet{+}{writeStructure(plant: Plant) : Promise$<$boolean$>$}{Persiste la struttura di un impianto nella cache}
\makeItf{WriteCachePort}{Porta di uscita per persistere la struttura di un impianto in cache}

% ── ENTITY ───────────────────────────────────────────────────────────────────

% infrastructure\persistence\entities\room.entity.ts
\addMet{+}{toDomain(entity: RoomEntity) : Room}{Converte l'entità persistita nel corrispondente oggetto di dominio}
\addMet{+}{fromDomain(room: Room) : RoomEntity}{Converte l'oggetto di dominio nella corrispondente entità persistita}
\makeCl{RoomEntity}{Entità di persistenza di una stanza, con metodi di conversione da/verso il modello di dominio}

% ── ADAPTER ──────────────────────────────────────────────────────────────────

% adapters\out\fetch-cache.adapter.ts
\addAtr{-}{getValidTokenPort: GetValidTokenPort}{Porta per ottenere un token di autenticazione valido}
\addAtr{-}{fetchNewCacheRepo: FetchNewCacheRepoPort}{Repository per recuperare la struttura di un impianto via HTTP}
\addMet{+}{fetch(cmd: FetchNewCacheCmd) : Promise$<$Plant$>$}{Ottiene un token valido e delega il fetch della struttura al repository}
\makeCl{FetchNewCacheAdapter}{Adapter che implementa FetchNewCachePort coordinando autenticazione e recupero dati}

% adapters\out\get-all-plantids.adapter.ts
\addAtr{-}{repo: GetAllPlantIdsRepoPort}{Repository per recuperare gli ID impianto via HTTP}
\addAtr{-}{getValidTokenPort: GetValidTokenPort}{Porta per ottenere un token di autenticazione valido}
\addMet{+}{getAllPlantIds() : Promise$<$string[]$>$}{Ottiene un token valido e delega la lettura degli ID al repository}
\makeCl{GetAllPlantIdsAdapter}{Adapter che implementa GetAllPlantIdsPort coordinando autenticazione e recupero degli ID}

% adapters\out\write-cache.adapter.ts
\addAtr{-}{writeCacheRepo: WriteCacheRepoPort}{Repository per persistere la struttura di un impianto}
\addMet{+}{writeStructure(plant: Plant) : Promise$<$boolean$>$}{Delega la persistenza della struttura dell'impianto al repository}
\makeCl{WriteCacheAdapter}{Adapter che implementa WriteCachePort delegando la scrittura al repository di persistenza}

% ── REPOSITORY (contratti) ────────────────────────────────────────────────────

% application\repository\fetch-new-cache.repository.ts
\addMet{+}{fetch(validToken: string, plantId: string) : Promise$<$PlantDto $|$ null$>$}{Recupera la struttura di un impianto dall'API esterna dato token e plantId}
\makeItf{FetchNewCacheRepoPort}{Contratto del repository HTTP per recuperare la struttura di un impianto}

% application\repository\get-all-plantids.repository.ts
\addMet{+}{getAllPlantIds(validToken: string) : Promise$<$string[]$>$}{Recupera la lista di tutti gli ID impianto dall'API esterna dato un token valido}
\makeItf{GetAllPlantIdsRepoPort}{Contratto del repository HTTP per ottenere tutti gli ID impianto}

% application\repository\write-cache.repository.ts
\addMet{+}{write(plant: PlantEntity) : Promise$<$boolean$>$}{Persiste una PlantEntity nella cache}
\makeItf{WriteCacheRepoPort}{Contratto del repository di persistenza per salvare la struttura di un impianto}

% ── IMPLEMENTAZIONE DEI REPOSITORY ───────────────────────────────────────────

% infrastructure\http\fetch-plant-structure-impl.ts
\addAtr{-}{httpService: HttpService}{Client HTTP per le chiamate alle API esterne}
\addMet{+}{fetch(validToken: string, plantId: string) : Promise$<$PlantDto $|$ null$>$}{Recupera la struttura completa di un impianto orchestrando le chiamate a stanze, dispositivi e datapoint}
\addMet{+}{getAllPlantIds(validToken: string) : Promise$<$string[]$>$}{Recupera la lista di tutti gli ID impianto dall'API esterna}
\addMet{-}{fetchRoom(validToken: string, plantId: string, room: ApiRoomDto) : Promise$<$RoomDto$>$}{Recupera i dettagli di una singola stanza e i suoi dispositivi}
\addMet{-}{fetchDevice(validToken: string, plantId: string, device: any) : Promise$<$DeviceDto$>$}{Recupera i dettagli di un singolo dispositivo e i suoi datapoint}
\addMet{-}{fetchDatapoint(datapoint: ApiDatapointDto) : DatapointDto}{Mappa un datapoint grezzo dell'API nel DTO interno}
\makeCl{FetchStructureCacheImpl}{Implementazione HTTP che recupera la struttura completa di un impianto componendo chiamate per stanze, dispositivi e datapoint}

% infrastructure\persistence\structure-cache-repository-impl.ts
\addMet{+}{write(plant: PlantEntity) : Promise$<$boolean$>$}{Salva la struttura dell'impianto nel sistema di persistenza della cache}
\makeCl{StructureCacheImpl}{Implementazione del repository che persiste la struttura di un impianto nella cache}